
%%%%%%%%%%%%%%%%%%%%%%%%%%%%%%%%%%%%%%%%%%%%%%%%%%%%%%%%%%%%%
\subsubsection{模型建立}

% ==== 开篇引文 ====
% 示例（A092风格）：问题一要求我们求解给定条件下XXX的XXX、XXX以及XXX。首先，我们可以通过...
针对问题一，本节将从...角度建立...模型。具体过程如下。

% ==== 步骤1：<标题>（A092风格用编号小节 "（1）标题" 加粗）====
\textbf{步骤1：...}

% 第1句：引入变量和场景
首先，...。设...为...，则有：
\begin{equation}
    ... = ...
\end{equation}
其中，$...$表示...，$...$表示...，$...$表示...。

将...代入...，可得：
\begin{equation}
    ... = ...
\end{equation}
进一步展开得：
\begin{equation}
    ... = ...
\end{equation}

% ==== 步骤2：<标题> ====
\textbf{步骤2：...}

根据...的性质，可以建立以下关系：
\begin{equation}
    ... = ...
\end{equation}
其中，$...$为...，$...$为...。

由式(1)和式(2)联立可得：
\begin{equation}
    ... = ...
\end{equation}

% ==== 步骤3：<标题> ====
\textbf{步骤3：...}

% ...（按需增加步骤，每步≥6个公式）...

% ==== 步骤4：求解算法 ====
\textbf{步骤4：...算法}

% 算法流程公式 + 参数表

\begin{longtable}{>{\centering\arraybackslash}p{0.16\textwidth}>{\centering\arraybackslash}p{0.24\textwidth}>{\centering\arraybackslash}p{0.24\textwidth}>{\centering\arraybackslash}p{0.24\textwidth}}
  \caption{算法参数设置}
  \label{tab:参数1} \\
  \toprule
  参数名 & 参数含义 & 取值 & 选取依据 \\
  \midrule
  \endfirsthead
  \caption{算法参数设置（续）} \\
  \toprule
  参数名 & 参数含义 & 取值 & 选取依据 \\
  \midrule
  \endhead
  \bottomrule
  \endfoot
  ... & ... & ... & ... \\
\end{longtable}

% ==== 总结桥接 ====
综上所述，本节完成了...模型的建立，得到了...，接下来将基于此进行求解。

% ============================================================
\subsubsection{模型求解}

% ==== 引文（A092风格：直接陈述求解结果，含数据来源）====
% 示例：基于上述模型，本节利用...数据进行求解。
基于上述模型，本节利用...数据进行求解。

% ==== 结果表格 ====
求解得到以下结果：

\begin{longtable}{>{\centering\arraybackslash}p{0.16\textwidth}>{\centering\arraybackslash}p{0.24\textwidth}>{\centering\arraybackslash}p{0.24\textwidth}>{\centering\arraybackslash}p{0.24\textwidth}}
  \caption{...结果}
  \label{tab:结果1} \\
  \toprule
  指标1 & 指标2 & 指标3 & 指标4 \\
  \midrule
  \endfirsthead
  \caption{...结果（续）} \\
  \toprule
  指标1 & 指标2 & 指标3 & 指标4 \\
  \midrule
  \endhead
  \bottomrule
  \endfoot
  ... & ... & ... & ... \\
\end{longtable}

% ==== 结果图 ====
如图\ref{fig:result1}所示，...（一句话点出图的结论）。

\begin{figure}[H]
  \centering
%   \includegraphics[width=0.8\textwidth]{../求解/问题一/图片/xxx.png}  % 求解后替换实际路径
  \caption{...}
  \label{fig:result1}
\end{figure}

% ==== 结果分析 ====
由表\ref{tab:结果1}可知，...。从图\ref{fig:result1}可以看出，...。综合以上结果，...。

% ==== 模型检验（A092在每问求解末尾嵌入检验，非单独章节）====
\paragraph{模型检验}

% 随机生成可行解进行比较（A092的验证方式）
% 示例：我们在满足约束条件的情况下，随机生成了100组可行解，计算其XXX，与本问中求得的最优XXX进行比较，如图X所示。我们发现...有力地说明了我们的...是有效且合理的。
为了验证模型的有效性，我们在满足约束条件的情况下，随机生成了100组可行解，计算其...，与本问中求得的最优...进行比较，如图\ref{fig:check1}所示。

\begin{figure}[H]
  \centering
%   \includegraphics[width=0.8\textwidth]{../求解/问题一/图片/模型检验.png}  % 求解后替换实际路径
  \caption{随机生成的100组可行解与最优解的比较}
  \label{fig:check1}
\end{figure}

我们发现...，这有力地说明了我们在问题一中构建的模型是有效且合理的。
